\chapter{Female Orgasm}
\index{Female Orgasm}

\section*{Definition and Overview}
The female orgasm is a peak of intense pleasure and release of sexual tension, accompanied by rhythmic contractions of the pelvic floor muscles and a wave of physical and emotional sensation. It is a normal, healthy part of sexuality, but it is also highly individual: how women reach orgasm, how easily, and how often varies enormously. Many women do not reach orgasm from penetrative intercourse alone, which is entirely normal. When difficulty reaching orgasm causes distress, it is called anorgasmia or delayed orgasm, and it is one of the most common sexual concerns women bring to health professionals. This chapter explains how female orgasm works, why it can be difficult to achieve, and what can help.

\section*{Epidemiology}
Surveys consistently show that while most women can reach orgasm, many do so inconsistently and a substantial minority rarely or never do. Around 70--80\% of women report that direct clitoral stimulation is the most reliable route to orgasm, while only 20--30\% reach orgasm from intercourse alone. Approximately 10--15\% of women have never experienced an orgasm, and many more experience situational difficulties at some stage of life. Distress about orgasm difficulties is common and is one of the most frequent reasons for seeking sexual health advice. Rates of satisfaction rise with better understanding, communication, and targeted help.

\section*{Aetiology and Pathophysiology}
\begin{itemize}
    \item \textbf{Normal physiology:} the clitoris, with its rich nerve supply, is the primary organ of female sexual pleasure; stimulation of the clitoris, vulva, and other erogenous areas builds arousal towards orgasm
    \item \textbf{Blood flow and arousal:} arousal increases blood flow to the genital tissues; orgasm involves a build-up of tension and its release through pelvic muscle contractions
    \item \textbf{Psychological factors:} anxiety about performance, body image concerns, stress, distraction, and past trauma inhibit the relaxation needed for orgasm
    \item \textbf{Relationship factors:} poor communication, mismatched expectations, and lack of partner involvement in clitoral stimulation
    \item \textbf{Medical conditions:} diabetes, neurological disease, pelvic surgery (including hysterectomy and episiotomy complications), and pelvic floor dysfunction
    \item \textbf{Medicines:} antidepressants, particularly SSRIs, commonly delay or abolish orgasm
    \item \textbf{Hormones:} low oestrogen (vaginal dryness and altered sensation) and low testosterone can reduce arousal and sensitivity
    \item \textbf{Alcohol and drugs:} depress sexual response and sensitivity
    \item \textbf{Lack of knowledge:} many women and couples do not know that clitoral stimulation is central to female orgasm
\end{itemize}

\section*{Clinical Features}
\begin{itemize}
    \item Inability to reach orgasm despite adequate stimulation and arousal, in situations where the woman wants to
    \item Delayed orgasm, taking much longer than expected or than previously
    \item Orgasms only in specific situations, such as alone but not with a partner
    \item Lifelong (never experienced) versus acquired (previously experienced) patterns
    \item Associated features that suggest causes: vaginal dryness, pain, low mood, medication use, or relationship distress
    \item Distress, frustration, avoidance of sex, and effects on the relationship
    \item It is important to distinguish difficulty reaching orgasm from low desire, which is a separate problem
\end{itemize}

\section*{Differential Diagnosis}
\begin{itemize}
    \item Female orgasmic disorder (anorgasmia), the specific diagnosis when delay or absence causes distress
    \item Low sexual desire, which reduces arousal and makes orgasm harder
    \item Sexual pain disorders such as dyspareunia and vaginismus
    \item Genitourinary syndrome of menopause with dryness and reduced sensation
    \item Medication-induced sexual dysfunction, especially with antidepressants
    \item Depression, anxiety, and the effects of trauma
    \item Neurological and endocrine conditions, including diabetes and thyroid disease
\end{itemize}

\section*{When to Seek Medical Care}
Women who are distressed by difficulty reaching orgasm, or who notice a change after starting a medicine or a health event, should discuss it with a doctor, psychologist, or sexual health service. It is a legitimate and common reason for help.

\textbf{Contact your local emergency services for:}
\begin{itemize}
    \item Sudden loss of sensation or function with neurological symptoms such as weakness, numbness, or headache, which requires urgent assessment
    \item Any emergency health event during or after intercourse, such as chest pain or collapse
    \item Suicidal thoughts, which can accompany underlying depression
\end{itemize}

\section*{Diagnosis and Investigations}
\begin{itemize}
    \item \textbf{History:} lifelong or acquired, situations in which orgasm is or is not achieved, medications, health conditions, relationship context, and levels of distress
    \item \textbf{Physical examination:} where indicated, including pelvic examination if pain, dryness, or sensation changes are present
    \item \textbf{Medication review:} identifying antidepressants and other drugs that delay orgasm
    \item \textbf{Blood tests:} for diabetes, thyroid disease, and hormone levels when a physical cause is suspected
    \item \textbf{Psychological assessment:} exploring anxiety, mood, body image, and trauma
    \item \textbf{No routine investigations:} most cases are understood through the history alone
\end{itemize}

\section*{Management}
\begin{itemize}
    \item \textbf{Education:} understanding that clitoral stimulation is the most reliable route to orgasm and that intercourse alone is not the norm
    \item \textbf{Sensate focus and solo exploration:} masturbation practice and structured exercises that build comfort with one's own body and response
    \item \textbf{Partner communication:} guiding partners on what works, and exploring mutual stimulation
    \item \textbf{Address contributing factors:} lubricants for dryness, pain management, stress reduction, and sleep
    \item \textbf{Medication adjustment:} with a doctor, reducing the dose, switching, or adding strategies to counter antidepressant-related delay
    \item \textbf{Sex therapy and CBT:} effective for anxiety-related and lifelong difficulties
    \item \textbf{Pelvic floor physiotherapy:} for women with pelvic floor dysfunction or after surgery
    \item \textbf{Vibrators and toys:} many women find external vibrators help substantially
    \item \textbf{Device and medical options:} clitoral vacuum devices and specialist-prescribed treatments are options in selected cases
\end{itemize}

\section*{Complications}
\begin{itemize}
    \item Distress, frustration, and reduced self-esteem
    \item Avoidance of sex and strain on the relationship
    \item Anxiety that worsens the problem, creating a vicious cycle
    \item Missed diagnosis of underlying conditions such as diabetes or pelvic floor issues
    \item Reduced treatment of the condition if the issue is never raised
\end{itemize}

\section*{Prognosis and Outcomes}
The outlook is generally very good. Most women with orgasm difficulties improve with education, communication, and targeted practice, and sex therapy helps even lifelong difficulties in most cases. Medication-related problems resolve when the medicine is adjusted. Understanding that orgasm from intercourse alone is not the norm relieves pressure for many women and couples, often improving sexual satisfaction even without further treatment. Persisting difficulties respond to specialist sexual health care, which is effective and widely available.

\section*{Prevention and Self-Care}
\begin{itemize}
    \item Learn about your own body and what stimulation works for you
    \item Share that knowledge openly with a partner
    \item Make time for unhurried intimacy, including foreplay and clitoral stimulation
    \item Reduce distractions and performance pressure during sex
    \item Manage stress, sleep well, and limit alcohol before sex
    \item Use lubrication if dryness is an issue
    \item Ask a doctor about sexual side effects when starting a new medicine
    \item Seek help early; sexual problems are common and treatable
\end{itemize}

\section*{Sources and Further Reading}
The following reputable sources provide further information for healthcare students and patients seeking reliable, current advice about female orgasm.
\begin{itemize}
    \item Healthdirect Australia. \textit{Female sexual problems.} https://www.healthdirect.gov.au/
    \item Jean Hailes for Women's Health. \textit{Orgasms and sexual pleasure.} https://www.jeanhailes.org.au/
    \item True Relationships \& Reproductive Health. \textit{Sexual health information.} https://www.true.org.au/
    \item NHS. \textit{Female sexual problems.} https://www.nhs.uk/conditions/sexual-problems-in-women/
    \item Mayo Clinic. \textit{Female sexual dysfunction.} https://www.mayoclinic.org/
    \item MedlinePlus. \textit{Female sexual dysfunction.} https://medlineplus.gov/
\end{itemize}
